\subsection{A coherent pure-complex lift of the rank-three kernels}

The coefficient-object alternative has a further constructive
refinement. The actual relation maps of \fid{F56} can be carried
by pure complexes, provided the relation between formal grading,
cohomological degree and Frobenius is specified. The full-Tate
and half-step conventions have different domains.

\begin{finding}[\fid{F61}: pure graded kernels with actual relation maps]
\label{f:purerankthreekernels}
Fix a prime power \(q\), \(\ell\nmid q\), and a finite extension
\(E/\mathbb Q_\ell\), enlarged for the stated arithmetic
coefficients. Multiplicity spaces have trivial Frobenius.
Retain the actual
64-state kernels \(K_1,K_2\) of \fid{F56}, with their explicit
quadratic and degree-five isomorphisms.

In the full-Tate convention \(v=q^n\), the coefficient functor
\begin{equation}\label{eq:fullTatekernelLift}
 \Phi_{\mathrm T}(U)
   =\bigoplus_e U_e\otimes E(-e)[-2e]
\end{equation}
is additive symmetric strong monoidal. Applied entrywise, it
lifts the whole graded word category, with every grading shift,
to actual kernels of bounded pure point complexes of weight zero.
It preserves the specified relation maps and their coherence
comparisons, and its Euler traces are the original graded
characters evaluated at \(v=q^n\).

For \(v=q^{n/2}\), choose \(\lambda=\sqrt q>0\) under a specified
complex embedding and an \(\ell\)-adic realization in \(E\).
There are integral state potentials \(t_a\) for the same
actual kernels such that
\begin{equation}\label{eq:halfstepkernelLift}
 (\Phi_tK)_{ab}
   =\bigoplus_e (K_{ab})_e\otimes L_e[-e+t_b-t_a],
 \qquad F\mid L_e=\lambda^e,
\end{equation}
is a strong monoidal lift on the category of even shifts of
words. Every entry is pure of weight \(t_b-t_a\), and its Euler
trace is the original entry character at \(q^{n/2}\).
All closed-word trace objects are pure of weight zero.
The same statements hold for the faithful six-state version.
The formal entrywise graded bar also transports: it is ordinary
point duality in the full-Tate convention and the corrected
entrywise dual \(C_{ab}^\vee[2(t_b-t_a)]\) in the half-step
convention.

In odd characteristic, entrywise derived Gaussian transport
from \fid{F53} gives actual geometric derived-kernel maps with
these purity properties. Multidegree objects need not be
perverse in degree zero. In both conventions, pure complex
weight zero permits several Frobenius radii on total
cohomology. No common-radius assertion, pointwise/dual
estimate, Hard Lefschetz structure, or designated canonical
basis follows from this construction.
\end{finding}

\begin{proof}
\par\smallskip\noindent\textbf{Convention and the full-Tate functor.}\quad
Let \(\mathcal V\) be finite \(\mathbb Z\)-graded \(E\)-vector
spaces with ordinary tensor symmetry, without Koszul signs on
the formal grading. The kernels have convolution
\[
 (K\star L)_{ab}=\bigoplus_h K_{ah}\otimes L_{hb}.
\]
They act on row vectors, so \(T_L\circ T_K\simeq T_{K\star L}\).
Their associator preserves the state path and the ordered
homogeneous coefficient factors.

On the finite-field point, BBD's smooth/lisse criterion says
that a complex is pure of weight \(w\) precisely when its
\(H^i\) is punctually pure of weight \(w+i\).
Every constructible point sheaf is lisse
\cite[5.1.5, pp.~125--126; 5.1.8, pp.~126--127]{BBD}.
All objects here are explicit bounded split complexes,
initially with zero differential. Their stated maps are
chain maps, hence derived morphisms. We use this specified
image without claiming it contains all ambient derived
morphisms.

Geometric Frobenius on \(E(-e)\) is \(q^e\).
Its term in \eqref{eq:fullTatekernelLift} is in degree \(2e\)
and has weight \(2e\); hence each term and the finite sum are
pure of complex weight zero. Moreover,
\begin{equation}\label{eq:fullTatekernelTrace}
 \operatorname{Tr}_{\mathrm{Eul}}(F^n;\Phi_{\mathrm T}(U))
   =\sum_e\dim(U_e)q^{en}
   =\operatorname{ch}(U)(q^n).
\end{equation}
All cohomological degrees are even, so there are no Euler
minus signs.

On a homogeneous tensor use the standard identification
\(E(-e)\otimes E(-f)\simeq E(-(e+f))\), retaining the
multiplicity vectors in order. Shifts add as
\(-2e-2f=-2(e+f)\), giving the tensorator.
Every associativity route is the same map on the same
ordered pure tensor; the unit maps to \(E\) in degree zero.
The derived tensor symmetry contributes
\((-1)^{2e\cdot2f}=1\), so it agrees with the original
ordinary flip. This proves symmetric strong monoidality
on all graded coefficients, including every shift.

Entrywise application and finite distributivity give
\[
 \Phi_{\mathrm T}(K)\star\Phi_{\mathrm T}(L)
       \simeq\Phi_{\mathrm T}(K\star L).
\]
On each path this is the preceding tensorator.
The flattened-path associator, unit and all specified
tensor-context comparisons are therefore retained.
Every actual graded kernel map, its inverse when present,
and every equality among these maps is preserved.
Each word entry is pure of weight zero, as is the finite
sum of the diagonal entries defining its trace object.

\par\smallskip\noindent\textbf{The exact half-step parity condition.}\quad
Let \(L_e\) be the Weil line of
\eqref{eq:halfstepkernelLift}. Fix the tensor
identifications by tensor powers of \(L_1\) and its dual.
Every complex conjugate of \(\lambda^e\) has absolute
value \(q^{e/2}\); thus \(L_e\) is pure of weight \(e\).
It may also be chosen as a continuous point representation:
\(\lambda^e\) is an \(\ell\)-adic unit, and its cyclic
image extends to the procyclic Galois group. The point
representation convention is BBD, 5.1.11, pp.~127--128.
For even \(e=2j\) one can take \(L_{2j}=E(-j)\).
For odd \(e\) this is an arithmetic coefficient line,
with no geometric half-Tate-motive assertion.

For fixed integer potentials \(t_a\), restrict to kernels
with
\begin{equation}\label{eq:kernelparitycondition}
       (K_{ab})_e=0
       \quad\text{unless}\quad
       e\equiv t_b-t_a\pmod2.
\end{equation}
This full additive subcategory contains the unit.
It is closed under convolution because the two parity
conditions through \(h\) add to the condition from \(a\)
to \(b\). It is closed under even global grading shifts,
whereas an odd global shift generally leaves it.

In \eqref{eq:halfstepkernelLift}, the degree of the
cohomology term is
\[
                   i=e-t_b+t_a.
\]
It is even by \eqref{eq:kernelparitycondition}.
The term's cohomology weight is \(e\), so its complex
weight is \(e-i=t_b-t_a\). All summands of that entry have
the same complex weight. Thus the entry is pure of
weight \(t_b-t_a\), and
\begin{equation}\label{eq:halfstepkernelTrace}
 \operatorname{Tr}_{\mathrm{Eul}}
   (F^n;(\Phi_tK)_{ab})
     =\operatorname{ch}(K_{ab})(q^{n/2}).
\end{equation}
The parity condition is also necessary for this effective
entry and prescribed pure weight: by
Finding~\ref{f:purepointLaurentcriterion}, a positive
coefficient at \(\lambda^e\) can occur only in degree
\(e-t_b+t_a\), which must be even. This necessity is
specific to the stated entry-weight comparison.

The convolution tensorator is given on homogeneous paths
by \(L_e\otimes L_f\simeq L_{e+f}\) and the shift identity
\begin{equation}\label{eq:kernelpotentialtelescoping}
 (-e+t_h-t_a)+(-f+t_b-t_h)
       =-(e+f)+t_b-t_a.
\end{equation}
It retains the state path and ordered multiplicity tensor.
All associativity and unit comparisons therefore agree
on pure tensors. Every cohomological factor in these
images is even, so the specified coefficient permutations
introduce no Koszul signs. This proves the strong
monoidal kernel lift.
Along a word, the potentials telescope as in
\eqref{eq:kernelpotentialtelescoping}; entries have weight
\(t_b-t_a\), and closed-word trace objects have weight zero.

\par\smallskip\noindent\textbf{The actual states and relation maps.}\quad
Use the coefficient spaces of \fid{F56},
\[
 A_1=E\langle-1\rangle\oplus E\langle1\rangle,\quad
 A_2=E\langle-2\rangle\oplus E\langle2\rangle,\quad
 \mathcal D=A_1^{\otimes2}.
\]
On a two-state block \(N(A)\) the actual kernels are
\[
 K_1=\begin{pmatrix}\mathcal D&0\\A&0\end{pmatrix},
 \qquad
 K_2=\begin{pmatrix}0&A\\0&\mathcal D\end{pmatrix}.
\]
The 64 states comprise 20 scalar-plus states, 4
scalar-minus states, 16 \(N(A_1)\) blocks and 4
\(N(A_2)\) blocks. Set \(t=0\) on scalar states,
\(t=(0,1)\) on each \(N(A_1)\), and \(t=(0,0)\)
on each \(N(A_2)\). Diagonal \(\mathcal D\) has
even degrees; \(A_1\) has odd degrees and \(A_2\)
has even degrees. Both named kernels satisfy
\eqref{eq:kernelparitycondition}.
For example,
\[
 (\Phi_tK_1)_{21}=L_{-1}\oplus L_1[-2],
 \qquad
 (\Phi_tK_2)_{12}=L_{-1}[2]\oplus L_1
\]
on an \(N(A_1)\) block. These entries have weights
\(-1\) and \(+1\), with even cohomological degrees.
The same potentials give the faithful six-state version.

The precise algebra presentation being lifted is IV,
Theorem~11.16, equations~(111)--(112), p.~65
\cite{IV}. Its explicit quadratic and degree-five kernel
maps were proved in \fid{F56}; here their transport,
rather than just their decategorified identity, is checked.
The quadratic map preserves the initial edge and final
reset loop. In a repeated power it retains the
chronological list of \(\mathcal D\) factors.
If coefficients are moved to the left, the two triple
reductions differ by the ordinary transposition of two
loop factors. Their images have even degrees, so that
same transposition gives the comparison, without a sign.
All the ordered permutation identities are retained.

For the degree-five map put
\[
 U=A_1^{\otimes2},\quad V=A_2^{\otimes2},\quad
 S=U\oplus V,\quad T=U\otimes V.
\]
Writing \(V_1=K_{121}\), \(W_1=K_{12121}\) and
using mirrored \(V_2,W_2\), the actual isomorphism is
\begin{equation}\label{eq:puredegreefivemap}
 W_1\oplus V_2S\oplus K_1T
       \simeq W_2\oplus V_1S\oplus K_2T.
\end{equation}
For completeness, on an \(A_1\) block the four left and
right lists are
\[
 (K_1UU,K_2UU,K_2UV,K_1UV),\qquad
 (K_2UU,K_1UU,K_1UV,K_2UV).
\]
The map sends positions \(1,2,3,4\) to \(2,1,4,3\),
using coefficient identities. On an \(A_2\) block the
lists are
\[
 (K_1VV,K_2VU,K_2VV,K_1UV),\qquad
 (K_2VV,K_1VU,K_1VV,K_2UV).
\]
Send positions \(1,2,3,4\) to \(3,4,1,2\), using the
\(U,V\) flip on the second and fourth maps.
On plus states use termwise identities, and on minus
states the zero isomorphism.
The alternating-word comparisons retain chronological
edge factors.

Every \(U,V,S,T\) degree is even, so these are allowed
central coefficients in the half-step subcategory.
The isolated odd object \(A_1\) is not such a central
coefficient; only its permitted off-diagonal appearances
are used. Thus every intermediate object and map in
\eqref{eq:puredegreefivemap} lies in the stated domain.
Both lifts carry these maps and their inverses to actual
chain isomorphisms. Their repeated copies commute with
the multiplicity matrix units, since corresponding
copies have identical potentials and identical maps.
Functoriality retains compositions, tensor contexts,
the specified associator and symmetry comparisons,
and disjoint-context interchange.

This is coherence of the constructed category and its
actual relation maps. IV's algebra presentation does
not prescribe higher categorical relations among all
rewrite paths; no assertion that every unspecified
relation loop is the identity is being added.

\par\smallskip\noindent\textbf{The precise formal bar.}\quad
The F56 bar supplement uses entrywise graded duality,
without transposing the states, and the order-preserving
product pairing as its contravariant monoidal tensorator.
For the full-Tate functor,
\[
 \Phi_{\mathrm T}(\mathsf B K)_{ab}
     \simeq\Phi_{\mathrm T}(K)_{ab}^{\vee},
\]
since the dual of \(E(-e)[-2e]\) is \(E(e)[2e]\).
These identifications commute with the tensorator and
double-dual evaluation on every homogeneous pure tensor.

Write \(d_{ab}=t_b-t_a\). The ordinary dual of a half-step
summand \(U_e\otimes L_e[-e+d_{ab}]\) is
\(U_e^\vee\otimes L_{-e}[e-d_{ab}]\), while applying
\(\Phi_t\) to its graded dual gives
\(U_e^\vee\otimes L_{-e}[e+d_{ab}]\). Therefore
\begin{equation}\label{eq:purekernelbarcorrection}
 \Phi_t(\mathsf B K)_{ab}
      \simeq\Phi_t(K)_{ab}^{\vee}[2d_{ab}].
\end{equation}
Plain point duality is insufficient when \(d_{ab}\ne0\).

Define \((\mathsf B_t C)_{ab}=C_{ab}^\vee[2d_{ab}]\)
on the even split kernel image. Finite duality commutes
with the intermediate-state sum and
\(d_{ah}+d_{hb}=d_{ab}\). Hence the product pairing gives
a tensorator, in each path summand,
\[
 C_{ah}^\vee[2d_{ah}]
       \otimes C'_{hb}{}^\vee[2d_{hb}]
   \simeq (C_{ah}\otimes C'_{hb})^\vee[2d_{ab}].
\]
Every shift and cohomological degree here is even.
No Koszul sign changes the product pairing.
For a triple tensor, both coherence routes evaluate
to the same product of three scalar pairings in the
same state summand. The diagonal unit has \(d_{aa}=0\);
its constraint is ordinary scalar evaluation.
Naturality follows by precomposing the two factors
with their respective maps. Thus \(\mathsf B_t\)
is contravariant strong monoidal, preserving the
order of kernel multiplication.

Moreover,
\[
 \mathsf B_t^2(C)_{ab}
   =(C_{ab}^\vee[2d_{ab}])^\vee[2d_{ab}]
   \simeq C_{ab}^{\vee\vee}.
\]
The canonical evaluation is its coherent involution:
on pure tensors the double-dual tensorator is the
product evaluation, and applying the involution twice
returns the original functional. This checks the
monoidal and involutive identities.
If \(C_{ab}\) has weight \(d_{ab}\), its dual has
weight \(-d_{ab}\); the shift by \(2d_{ab}\)
restores weight \(d_{ab}\). The corrected bar thus
preserves the entry-weight specification.

The specified F56 self-dual forms and their inverse-map
identities transport through these comparisons.
Unshifted named kernels are self-dual; their explicit
relation maps are isometries and hence dualize to
their inverses under those identifications.
In general contexts, the contexts themselves are
replaced by their bars, with grading shifts reversed.
No self-duality of every shifted object is asserted.
On Euler characters this is \(v\mapsto v^{-1}\).
It is neither a positive Hermitian polarization nor
a transpose-dual action adjoint.

\par\smallskip\noindent\textbf{Derived Gaussian transport.}\quad
In odd characteristic use \fid{F53}:
\[
 L=\mathcal L_\psi(x^2),\quad
 G=H_c^1(\mathbb A^1_{\bar k},L),\quad
 J(C)=L\otimes\underline{C\otimes G^\vee}[1].
\]
The line \(G\) has weight one. Its geometric product is
\[
 A\star_\psi B=
 \underline{G(1)}[1]\otimes
 Rm_!\left(A\boxtimes B\otimes\mathcal L_\psi(2xy)\right),
 \qquad m(x,y)=x+y.
\]
Use the transported associator and symmetry from
\fid{F53}, and its image unit \(J(E)\).
In particular the commutor is not the naive swap of
two shifted Gaussian kernels.

For the explicit split point complexes define
\[
 \widetilde J\left(\bigoplus_i C_i[-i]\right)
       =\bigoplus_i J(C_i)[-i].
\]
Equivalently this is the same derived formula for \(J\).
Expand two finite sums and, on each \((i,j)\) summand,
use the natural F53 comparison with total shift
\([-i-j]\). This gives
\[
 \widetilde J(C)\star_\psi\widetilde J(C')
       \simeq\widetilde J(C\otimes C').
\]
The associativity and unit diagrams restrict to the
F53 diagrams with the same total shift.
Its derived transported symmetry includes the
coefficient Koszul sign; for our even point degrees
that extra sign is one. The existing Gaussian
symmetry correction is therefore preserved.

If \(C_i\) is pure of weight \(w+i\), then \(J(C_i)\)
has lisse coefficient of weight \(w+i-1\) and shift
\([1]\), so it is perverse pure of weight \(w+i\).
Consequently \(J(C_i)[-i]\) is a pure derived complex
of weight \(w\). Finite sums of this weight remain pure
by BBD's criterion.
Entrywise transport, finite intermediate-state sums
and \(\star_\psi\) therefore carry all the preceding
kernel maps with their stated weights.
On this split image the Gaussian compact-cohomology
evaluation is inverse to \(J\); it commutes with the
finite sums and shifts and recovers the point Euler
traces.
The shifted images need not be perverse in degree zero.
No global monoidality or exactness statement for all
constructible complexes is inferred.
The bar just proved transports along this specified
image equivalence and its compact-cohomology inverse.
It is not literal Verdier duality on the line:
\fid{F53} distinguishes that dual, the internal
coefficient dual, and the reciprocal weight system.
\end{proof}

\paragraph{What this supplies to a purity specification.}
Derived purity of the specified coefficient kernels
and their actual relation maps is now achieved at
rank three. It is different from a single Frobenius
radius on ungraded total cohomology. For example,
\[
 \Phi_{\mathrm T}(\mathcal D)
    =E(-2)[-4]\oplus E^{\oplus2}\oplus E(2)[4],
\]
whereas the half-step diagonal entry is
\[
       E(-1)[-2]\oplus E^{\oplus2}\oplus E(1)[2].
\]
Their total cohomology has radii \(q^2,1,q^{-2}\)
and \(q,1,q^{-1}\), respectively. Each degree \(i\)
has weight \(i\); that is the reason their complex
weight is zero. Word amplitudes grow with word length.
This is not a fixed finite Weil space supporting
all word traces.

The pointwise and independent dual estimates of
\fid{F50} are unchanged. No Hard Lefschetz operator,
polarization or designated canonical labels are
supplied. The full-Tate construction transports the
whole F56 category and its completion limitations;
the half-step construction transports the explicitly
restricted even-shift category. Neither turns it into
the source's desired canonical-simple category.

The actual full-algebra rank-four \(N^\#\) block of
\fid{F57} prevents extending this particular pure
word-trace comparison to every module and rank.
The half-step obstruction is
Finding~\ref{f:purepointLaurentcriterion}.
In the full-Tate convention all roots \(q^e\) have
even weights \(2e\), so a pure point complex of any
one weight would again force all visible coefficients
to have one parity sign. The same mixed-sign word
trace fails that condition.
This does not refute the printed signed/factored
positivity or exclude mixed geometric coefficients.
The designated arithmetic/grading/basis comparison
and general canonical positivity remain unproved.

% Separate paragraph for Part III.
\paragraph{\fid{F61} in Grothendieck's language.}
The source formal grading is an ordinary split-torus
character grading. SGA~3, I, 4.7.2--3, with cote~90,
scan~17, supplies the action framework; SGA~6, V,
\S\S2--3, with cote~133, scan~18, supplies representation
and \(\lambda/\sigma\) language for the symmetric
coefficient category \cite{SGA3,SGA6,Archive}.
No \(\lambda\)-ring structure is asserted on the
noncommutative word algebra.
The functors above prescribe extra grading-to-degree
and grading-to-Frobenius data; they are not consequences
of that dictionary.

Their bounded coefficient complexes use the perfect
framework of SGA~6, I, Definition~4.7, pp.~121--122,
and Example~4.8, p.~122 \cite{SGA6}.
BBD, 5.1.8, pp.~126--127, proves their exact
complex-weight claims \cite{BBD}.
The Gaussian transport uses the existing compact-trace
geometry of \fid{F53}: Milne, Proposition~18.4,
pp.~121--122; Theorem~22.1, Remarks~22.3(b),(d) and
Proposition~22.5, pp.~135--137; and
Theorem~24.1(a),(b), p.~144 \cite{Milne}.
This is a pure derived coefficient construction on the
specified image, with no archival canonical-basis or
positivity theorem added.
